%=====================================================================
% Implementación · Script de inserción de datos
%   Las cifras las genera bd/generar_datos_prueba.py. Los registros, lo
%   que representa cada uno y las capturas están en la sección Datos.
%=====================================================================
\subsection*{Script de inserción de datos}
\addcontentsline{toc}{subsection}{Script de inserción de datos}

\at{bd/insertar\_datos\_prueba.sql} carga dos cosas en una sola transacción: los catálogos que el sistema necesita para funcionar y un escenario de prueba de \dpNumFilas{} registros que ocupa las \dpNumTablasConDatos{} tablas. El escenario es una fotografía de la base el \dpFecha{} a las 18:00 UTC, con el juego en marcha: hay una partida en curso, una sala abierta, un jugador en la cola y una invitación pendiente. Si una fila falla, no se carga ninguna. Su listado completo está al final de esta sección; los registros de cada tabla, lo que representa cada cuenta y cada partida, y las capturas de la carga y de las pruebas están en la sección \textit{Datos}.

\subsubsection*{Catálogos}
\addcontentsline{toc}{subsubsection}{Catálogos}

Los catálogos toman de los casos de uso todo lo que estos fijan: las seis ranuras (\mbox{D-03}); los tres modos con sus tableros, plazas y muros (\mbox{CU-21} \mbox{RN-02}); los cuatro niveles de la IA con su fuerza (\mbox{CU-27} \mbox{RN-02}); las siete lecciones (\mbox{CU-41} \mbox{RN-01}); los cinco motivos de reporte (\mbox{CU-42} \mbox{RN-08}); y las divisiones bronce, plata, oro y muro de piedra de la tabla de dominios. Cada fila lleva su \at{codigo}, como \at{CLASICO} o \at{PEON\_CLASICO}, y no su nombre, que está en los diccionarios de recursos (\mbox{D-21}). Lo que ningún caso fija es de prueba y conviene confirmarlo: los códigos, los objetos y sus precios, los umbrales de las divisiones, los tipos de caja y sus probabilidades, y las palabras del filtro.

\subsubsection*{Coherencia con el modelo}
\addcontentsline{toc}{subsubsection}{Coherencia con el modelo}

El script no se escribió a mano. \at{bd/generar\_datos\_prueba.py} parte del escenario de \at{bd/datos\_prueba} (cuentas, partidas con sus jugadas, economía y comunidad) y deduce todo lo que el modelo guarda como consecuencia de otras filas, igual que lo harían los casos de uso:

\begin{reglas}
  \item \textbf{Las jugadas son legales.} Cada jugada se valida con las reglas del \mbox{CU-20} y el \mbox{CU-21}: pasos ortogonales, saltos al rival, muros de dos casillas que no se solapan ni se cruzan, turnos en orden y muros disponibles. La última jugada de una partida ganada por meta llega de verdad a la meta.
  \item \textbf{El estado del tablero sale de las jugadas.} La casilla, los muros restantes y el reloj de cada participación, y el turno de la partida, se calculan aplicándolas en orden. En una partida perdida por tiempo, el reloj del perdedor llega a cero en su turno, y la partida termina en ese momento (\mbox{CU-20} \mbox{FA-06}).
  \item \textbf{La progresión sale de las partidas.} El elo de entrada de cada partida es el de salida de la anterior; las estadísticas de cada modo, la división y su historial se acumulan partida a partida; y cada partida clasificatoria deja un movimiento de monedas por participante (\mbox{CU-25} \mbox{RN-08}).
  \item \textbf{El saldo sale de los movimientos.} Compras, recompensas, cajas y conversiones se aplican en orden de fecha; el saldo nunca queda negativo y el que se guarda es la suma de los movimientos (\mbox{CU-40} \mbox{RN-02}).
  \item \textbf{El idioma de cada aceptación sale de la cuenta.} Los términos se aceptan en el \at{idioma\_preferido} de cada cuenta (\mbox{D-21}), así que \at{sofia\_salto} los aceptó en inglés y las demás cuentas en español.
  \item \textbf{Lo que exige el alta.} Toda cuenta tiene los objetos iniciales, un objeto equipado por ranura y estadísticas en cada modo activo (\mbox{CU-02}); cada participación, el aspecto que su jugador llevaba al empezar (\mbox{CU-14} \mbox{RN-03}).
\end{reglas}

El script termina con \dpNumComprobaciones{} consultas que recalculan todo lo anterior en la propia base y cuentan las filas que no coinciden. La última compara en binario los \dpNumTextosUnicode{} textos con caracteres fuera de ASCII con su valor reconstruido carácter a carácter con \at{NCHAR}, de modo que no depende de cómo se leyó el archivo; en la ejecución de prueba, todas devuelven cero. Las contraseñas y los tokens de prueba se guardan como resúmenes calculados con \at{HASHBYTES}, sin valores en claro, aunque el servidor de partidas usará su propio algoritmo de derivación (\mbox{CU-01} \mbox{RN-02}).
